\documentclass[11pt,a4paper]{article}

\usepackage[T1]{fontenc}
\usepackage[utf8]{inputenc}


\usepackage{geometry}
\geometry{a4paper, margin=2.5cm}

\usepackage{amsmath}
\usepackage{amssymb}
\usepackage{amsthm}
\usepackage{booktabs}
\usepackage{tabularx}
\usepackage{array}
\usepackage{xcolor}
\usepackage{listings}
\usepackage{hyperref}
\usepackage{url}
\usepackage{graphicx}
\usepackage{enumitem}
\usepackage{abstract}
\usepackage{titlesec}
\usepackage{fancyhdr}
\usepackage{footnote}
\usepackage[english]{babel}
\usepackage{multirow}

\hypersetup{
  colorlinks=true,
  linkcolor=black,
  urlcolor=blue,
  citecolor=black,
  pdflang={en-US},
  pdfauthor={PrivateNess Network},
}

% Explicitly set PDF language tag for TTS / accessibility
\AtBeginDocument{\pdfcatalog{/Lang (en-US)}}

\definecolor{codeblue}{RGB}{30,100,180}
\definecolor{codegray}{RGB}{100,100,100}
\definecolor{codegreen}{RGB}{0,130,80}
\definecolor{backcolor}{RGB}{248,248,248}

\lstdefinestyle{ness}{
  backgroundcolor=\color{backcolor},
  commentstyle=\color{codegray}\itshape,
  keywordstyle=\color{codeblue}\bfseries,
  stringstyle=\color{codegreen},
  basicstyle=\ttfamily\small,
  breaklines=true,
  captionpos=b,
  keepspaces=true,
  showspaces=false,
  showstringspaces=false,
  frame=single,
  framesep=4pt,
  xleftmargin=8pt,
  xrightmargin=8pt,
}
\lstset{style=ness}

\theoremstyle{definition}
\newtheorem{definition}{Definition}[section]
\newtheorem{property}{Property}[section]

\pagestyle{fancy}
\fancyhf{}
\rhead{\small PrivateNess Network --- Sovereign Inference-as-a-Service}
\lhead{\small April 2026}
\cfoot{\thepage}

\title{\textbf{Sovereign Inference-as-a-Service}\\[0.4em]
  \large Decentralized AI Endpoint Discovery, Authentication,\\
  and Payment over Commodity Hardware}

\author{
  Jeff Bouchard\\
  PrivateNess Network\\
  \texttt{privateness.network}
}

\date{April 2026}

% ─────────────────────────────────────────────────────────────────────────────
\begin{document}

\maketitle
\thispagestyle{fancy}

% ─── ABSTRACT ────────────────────────────────────────────────────────────────
\begin{abstract}
Contemporary AI inference services rely on centralized infrastructure: corporate
API gateways, revocable authentication credentials, fiat payment rails, and
BGP-reachable endpoints subject to geofencing and regulatory seizure.  We present
a sovereign inference-as-a-service architecture that addresses these
dependencies without relying on any centralized operator.  Service discovery is performed via RFC~2782 SRV records
stored in the Emercoin Name-Value Storage (NVS) blockchain.  Trust is rooted
entirely within the NVS record chain: the \texttt{inf:} capability record
references the \texttt{node:} registration record, which in turn references
the operator's \emph{privatenessID} (\texttt{id:}) WORM~XML identity record
containing the authoritative public key.  No DNS record type is repurposed
for authentication; no X.509 hierarchy participates at any point.  Transport
is carried by the Skywire mesh network using public-key routing, eliminating
IP-layer visibility.  Payment is denominated in EMC micro-transactions, with
proof-of-stake staking rewards at a fixed rate of 6\% per annum producing a
net-negative operational cost model for sufficiently staked nodes:
the chain subsidizes node operators before they serve a single inference request.
Stochastic sampling draws from a sovereign UHEPRNG entropy source rather than
operator-controlled randomness, closing a subtle output-manipulation attack
surface.  The entire stack is validated on commodity hardware at multiple price
tiers---from Raspberry Pi~3 and above through directly-supported and
community-ported system-on-chip TV-box hardware---establishing a reproducible
minimum viable node floor at consumer electronics price points.  We formalize the
\texttt{ness:}~namespace registry, the WORM~XML schema hierarchy, and the
complete five-step discovery protocol, and identify open problems including a
provider reputation protocol and liveness signalling without record mutation.
\end{abstract}

\vspace{1em}
\noindent\textbf{Keywords:} decentralized inference, blockchain service discovery,
sovereign identity, mesh networking, commodity hardware, proof-of-stake economics,
entropy as a security primitive.

\newpage

% ─── 1. INTRODUCTION ─────────────────────────────────────────────────────────
\section{Introduction}

The dominant model of AI inference delivery in 2025 resembles any other
cloud service in its dependency structure: a corporate operator runs compute, sells access via
API keys, logs requests by default, and terminates service at will through Terms of
Service enforcement, payment processor pressure, or regulatory mandate.  Users
have no recourse, no alternative, and no visibility into what is done with their
prompts.

Decentralization efforts in adjacent domains---file storage \cite{ipfs},
computation \cite{ethereum}, communications \cite{skywire}---have demonstrated
that this dependency is architectural, not inevitable.  Yet the ``decentralized AI''
literature has largely addressed training data governance and model ownership while
leaving the inference layer---the point of actual use---in the same centralized
hands.

This paper addresses that gap.  We present a complete architecture for
\emph{sovereign inference-as-a-service}: AI inference that is discoverable without
a central registry, authenticated without a certificate authority, transported
without BGP visibility, paid without fiat rails, and operable on hardware costing
less than a restaurant meal.

The architecture is not speculative.  Every component is either deployed or
directly derivable from deployed infrastructure: Emercoin NVS has operated since
2013 \cite{emercoin}; the PrivatenessTools client library is in active use;
Skywire mesh nodes are running; the WORM XML format and privatenessID identity
system are anchored on-chain.  This paper formalizes the integration and
identifies the remaining open problems.

\subsection{Contributions}

\begin{enumerate}[leftmargin=*, label=\textbf{C\arabic*.}]
  \item \textbf{WORM-chained trust without DNS repurposing.}  We demonstrate
    that inference endpoint authentication can be rooted entirely within NVS
    WORM records---\texttt{inf:} $\to$ \texttt{node:} $\to$ \texttt{id:}---without
    repurposing any DNS record type for a purpose outside its specification.
    Trust terminates at the AuxPoW-hardened blockchain, not at a certificate
    authority, a DNS record, or any corporate infrastructure.

  \item \textbf{Proof-of-stake economic inversion.}  We demonstrate that the
    EMC fixed PoS rate of 6\% per annum can exceed NVS operation and renewal
    costs for a node operator holding a modest staking balance, producing a
    net-negative cost model.  Inference capacity is not a cost center for
    staking operators; the chain subsidizes it.

  \item \textbf{Standards-compliant sovereign service discovery.}  RFC~2782 SRV
    records stored in EmerDNS provide load balancing, failover, and multi-transport
    addressing using existing DNS client semantics, with no centralized registry
    dependency.

  \item \textbf{Entropy as a security primitive.}  We formalize the role of
    UHEPRNG-sourced sampling entropy in closing an operator-controlled output
    manipulation attack surface that is implicit in all centralized inference
    deployments but has not been previously articulated as a threat model.

  \item \textbf{Reproducible commodity hardware floor.}  We establish a
    multi-tier, costed hardware taxonomy---from single-board computers to
    community-ported system-on-chip devices---as a falsifiable claim about the
    minimum viable sovereign inference node.
\end{enumerate}

\subsection{Paper Structure}

Section~\ref{sec:background} reviews the component technologies.
Section~\ref{sec:architecture} presents the full stack.
Section~\ref{sec:namespace} defines the \texttt{ness:} namespace registry.
Section~\ref{sec:discovery} specifies the SRV/WORM discovery protocol.
Section~\ref{sec:economics} analyzes the staking economic model.
Section~\ref{sec:hardware} characterizes the hardware deployment tiers.
Section~\ref{sec:security} provides a threat model and security analysis.
Section~\ref{sec:limitations} states open problems.
Section~\ref{sec:conclusion} concludes.

% ─── 2. BACKGROUND ───────────────────────────────────────────────────────────
\section{Background}
\label{sec:background}

\subsection{Emercoin NVS and AuxPoW}

Emercoin is a proof-of-work blockchain launched in December 2013 as a fork of
Peercoin, with merged mining (AuxPoW) against Bitcoin, inheriting Bitcoin's
accumulated hash rate as a security subsidy \cite{emercoin}.  Its Name-Value
Storage subsystem provides a key-value store
with blockchain-level immutability: records are cryptographically committed to
the chain at write time and are not modifiable without the controlling private
key.  Expiry is block-count-based; records can be renewed but not silently
altered.  We treat NVS as a \emph{write-once-read-many} (WORM) store for all
structured PrivateNess records.

The security cost of rewriting NVS history is bounded below by the Bitcoin
network's accumulated hash rate, scaled by the AuxPoW coefficient defined by
Khovayko's hack51 framework \cite{hack51}.  This makes history revision
economically prohibitive for any adversary without a majority of global Bitcoin
hash rate.

\subsection{WORM XML and PrivatenessTools}

The WORM (World Object Markup Language) format is a structured XML schema
developed by Sokolov \cite{privatenesstools} for signed, chain-committed
records in the PrivateNess ecosystem.  Every WORM document carries a
\texttt{<schema>} URI declaring its type and version, an ownership reference
to a sovereign identity, and a creation timestamp.  The PrivatenessTools
library provides \texttt{keygen} for key material generation and \texttt{key
worm} for emitting signed WORM documents from a key file.

\subsection{privatenessID}

A privatenessID is a sovereign identity anchored in Emercoin NVS under the
\texttt{id:} prefix.  The NVS value is a WORM XML document (schema
\texttt{ness:worm:id:1.0}) containing the identity's public key, creation
block reference, and optional transport addresses.  The PrivateNess Network's
own identity is anchored at Emercoin block~\#445669 (November~1, 2020), establishing
a four-year prior art timestamp for the identity system.

\subsection{Skywire Mesh Network}

Skywire is a public-key-routed mesh network developed by Smietana and Narrow
\cite{skywire}.  Routing is performed on 33-byte compressed elliptic curve
public keys rather than IP addresses.  Traffic between nodes is not visible at
the IP layer; an observer with access to a Skywire node sees encrypted packets
addressed to public keys, not to BGP prefixes.  Bandwidth is accounted in Coin
Hours, Skycoin's native time-weighted currency unit.  A Skywire-only inference
node has no TCP/IP listening socket; it presents no surface to conventional
network scanning, geofencing, or ASN-level blocking.

\subsection{EmerDNS}

EmerDNS maps conventional DNS record types (A, AAAA, SRV, MX, TXT) to
Emercoin NVS records under the \texttt{dns:} prefix.  A DNS-compatible resolver
reads NVS records and serves them to standard DNS clients.  Because the records
are stored on-chain, they are resistant to seizure, redirection, or censorship by a
registrar or registry operator.  The \texttt{.lib}, \texttt{.bazar}, and
\texttt{.coin} TLDs are native; \texttt{.ness} is the PrivateNess network TLD.

\subsection{RFC~2782 SRV Records}

RFC~2782 \cite{rfc2782} defines the DNS SRV resource record type for
service location.  An SRV record of the form \texttt{\_service.\_proto.name}
carries priority, weight, port, and target fields.  Clients select among
multiple SRV records using a weighted random algorithm that provides
load balancing within a priority class and automatic failover to lower-priority
records.  SRV is supported natively by all major DNS client libraries.

% ─── 3. ARCHITECTURE ─────────────────────────────────────────────────────────
\section{System Architecture}
\label{sec:architecture}

Figure~\ref{fig:layers} shows the full stack.  We describe each layer in
ascending order.

\begin{figure}[ht]
\begin{center}
\begin{tabular}{|c|}
\hline
\rule{0pt}{2.2ex}\textbf{Application layer}: inference client / payment \\
\hline
\rule{0pt}{2.2ex}\textbf{Payment layer}: EMC micro-transactions per token \\
\hline
\rule{0pt}{2.2ex}\textbf{Inference layer}: llama.cpp / Ollama, UHEPRNG entropy \\
\hline
\rule{0pt}{2.2ex}\textbf{Transport layer}: Skywire mesh (pubkey routing) or inet \\
\hline
\rule{0pt}{2.2ex}\textbf{Discovery layer}: EmerDNS SRV + NVS WORM trust chain \\
\hline
\rule{0pt}{2.2ex}\textbf{Identity layer}: privatenessID WORM XML (\texttt{id:}) \\
\hline
\rule{0pt}{2.2ex}\textbf{Consensus layer}: Emercoin NVS + Bitcoin AuxPoW \\
\hline
\end{tabular}
\end{center}
\caption{PrivateNess sovereign inference stack (ascending).}
\label{fig:layers}
\end{figure}

\subsection{Consensus Layer}

All persistent state---identity records, node registrations, service discovery
records, capability metadata, and the namespace schema index---is stored in
Emercoin NVS.  No server, database, or CDN is required.  An adversary wishing
to remove or alter any record must control a majority of the Bitcoin network's
hash rate for the duration of the attack.  Records that have been on-chain for
years are practically immutable.

\subsection{Identity Layer}

Every actor in the system---node operators, inference clients, reputation
attestors---is identified by a privatenessID: an \texttt{id:} NVS record
containing a WORM XML document with the actor's public key.  The key algorithm
is Ed25519 for current deployments, with a migration path to ML-DSA-65 (FIPS
204) specified in the Identity Bedrock paper \cite{identitybedrock}.

The privatenessID is not an email address, a phone number, or a domain name.
It is a chain-committed cryptographic key binding.  No personal information
is required or implied.

\subsection{Discovery Layer}

Service discovery uses two mechanisms:
\begin{itemize}
  \item \textbf{SRV record}: locates the endpoint (protocol, host, port).
  \item \textbf{WORM trust chain}: authenticates the endpoint operator via
    the \texttt{inf:} $\to$ \texttt{node:} $\to$ \texttt{id:} record sequence.
\end{itemize}

The complete discovery protocol is specified in Section~\ref{sec:discovery}.

\subsection{Transport Layer}

Nodes declare one or more transport networks in their registration record:
\texttt{inet} (clearnet TCP/IP), \texttt{skywire}, \texttt{yggdrasil},
\texttt{i2p}, or \texttt{tor}.  Clients select the transport most consistent
with their threat model.  A Skywire-only node has no conventional IP address; it is not reachable via the standard internet.  For Skywire nodes, the SRV target is the node's
compressed public key suffixed with \texttt{.skywire}; port is set to 0 by
convention, as Skywire is transport-native and does not use TCP port numbers.

\subsection{Inference Layer}

The inference runtime (llama.cpp, Ollama, or equivalent) runs on the node's
local hardware.  The critical property is the entropy source for sampling:
temperature-based token selection requires a random number generator; on a
sovereign node, this is provided by the PrivateNess UHEPRNG service
\cite{pyuheprng} rather than the host operating system's CSPRNG seeded by
vendor-controlled sources.  The adversarial significance of this is formalized
in Section~\ref{sec:entropy-threat}.

\subsection{Payment Layer}

Inference is billed in EMC per 1,000 tokens.  EMC is a transferable,
on-chain asset with no fiat intermediary.  The pricing is declared in the
node's \texttt{inf:} WORM XML capability record (Section~\ref{sec:inf-prefix})
and is verifiable before a session is opened.  Billing granularity is
configurable; the WORM record carries a base unit of EMC/1k tokens.

% ─── 4. NAMESPACE REGISTRY ───────────────────────────────────────────────────
\section{The \texttt{ness:} Namespace Registry}
\label{sec:namespace}

All PrivateNess NVS records are subordinate to the \texttt{ness:} root namespace.
Table~\ref{tab:prefixes} lists the canonical prefixes.  WORM XML schemas are
themselves registered as \texttt{ness:schema:\textit{type}:\textit{version}} NVS
records, making the schema registry self-describing and chain-hosted.

\begin{table}[ht]
\centering
\caption{Canonical \texttt{ness:} namespace prefixes.}
\label{tab:prefixes}
\begin{tabularx}{\textwidth}{@{}llX@{}}
\toprule
\textbf{Prefix} & \textbf{WORM schema} & \textbf{Role} \\
\midrule
\texttt{id:} & \texttt{ness:worm:id:1.0} & Sovereign identity (privatenessID). Trust anchor. \\
\texttt{inf:} & \texttt{ness:worm:inf:1.0} & Inference node capability and pricing. \\
\texttt{node:} & \texttt{ness:worm:node:1.0} & General node transport addresses and services. \\
\texttt{dns:} & DNS RR (plain text) & EmerDNS records: A, AAAA, SRV, TXT. Transport addressing only; not used for authentication. \\
\texttt{meme:} & \texttt{ness:worm:meme:1.0} & Immutable content records (NESSMeme). \\
\texttt{rep:} & \texttt{ness:worm:rep:1.0} & Node reputation attestations (specified; not yet implemented). \\
\texttt{ness:schema:} & --- & Schema registry: canonical WORM schema definitions. \\
\texttt{ness:index:} & JSON array & Append-only enumeration indexes per type. \\
\texttt{ness:version} & plain text & Namespace specification version. \\
\bottomrule
\end{tabularx}
\end{table}

\subsection{Append-only Enumeration}
\label{sec:enumeration}

Standard DNS does not support wildcard zone enumeration; a client cannot
issue a query equivalent to ``list all SRV records under
\texttt{\_inference.\_tcp.*.private.ness}.''  EmerDNS NVS supports prefix
scans on the key-value store, but this requires NVS API access rather than
standard DNS queries.

We address this with \texttt{ness:index:\textit{type}} records: append-only
JSON arrays of known record keys.  A node registering an inference service
appends its \texttt{inf:node\_\textit{hash}} key to \texttt{ness:index:inf}.
Clients without direct NVS access use this index for discovery enumeration.
The append-only property is enforced by convention: index records are never
reduced.  Nodes that go offline remain in the index; liveness is determined
by transport-layer connectivity, not by index membership.

\subsection{The \texttt{id:} Prefix}
\label{sec:id-prefix}

\begin{definition}[privatenessID]
A \emph{privatenessID} is an Emercoin NVS record with key \texttt{id:\textit{username}}
whose value is a WORM XML document conforming to schema \texttt{ness:worm:id:1.0},
containing at minimum: a public key, a creation timestamp, and the Emercoin block
number at which the record was first committed.
\end{definition}

The block number serves as a tamper-evident timestamp.  All downstream records
(\texttt{inf:}, \texttt{node:}, \texttt{rep:}) reference their owner's
\texttt{id:} key, making every record in the namespace traceable to a sovereign
identity without disclosing personal information.

\subsection{The \texttt{inf:} Prefix}
\label{sec:inf-prefix}

An inference node registers a capability record with key
\texttt{inf:node\_\textit{H}} where \textit{H} is the first 16 hexadecimal
characters of the SHA-256 hash of the owning identity's public key.  The WORM
XML value (schema \texttt{ness:worm:inf:1.0}) carries:

\begin{itemize}
  \item \texttt{<author>}: the owning \texttt{id:} NVS key.
  \item \texttt{<node>}: the associated \texttt{node:} NVS key.
  \item \texttt{<model>}: the model identifier (GGUF tag or equivalent).
  \item \texttt{<context\_window>}: maximum context in tokens.
  \item \texttt{<price\_per\_1k>}: price in EMC per 1,000 tokens.
  \item \texttt{<endpoint>}: the HTTP path for the inference API.
  \item \texttt{<skywire\_pubkey>}: Skywire transport address (optional).
\end{itemize}

This record is advisory metadata.  The authoritative endpoint is determined by
the SRV record (Section~\ref{sec:discovery}).

\subsection{The \texttt{dns:} Prefix}

EmerDNS NVS records under \texttt{dns:} carry plain-text DNS resource record
values.  The key is the DNS name with \texttt{dns:} prepended.  Three record
types are relevant to inference discovery:

\begin{enumerate}
  \item \texttt{dns:\_inference.\_tcp.\textit{host}}: SRV record for TCP transport.
  \item \texttt{dns:\_inference.\_skywire.\textit{host}}: SRV record for Skywire transport (port~0).
  \item \texttt{dns:\textit{host}}: A or AAAA address record for inet transport targets.
\end{enumerate}

% ─── 5. DISCOVERY PROTOCOL ───────────────────────────────────────────────────
\section{Discovery and Authentication Protocol}
\label{sec:discovery}

\subsection{SRV Record Structure}

An inference node operator writes the following NVS record:

\begin{lstlisting}[caption={SRV record for TCP inference endpoint.}]
KEY:   dns:_inference._tcp.batman1.private.ness
VALUE: SRV 10 100 8080 batman1.private.ness
\end{lstlisting}

The four fields follow RFC~2782: \emph{priority}, \emph{weight}, \emph{port},
and \emph{target}.  Multiple nodes registering under the same service label
(\texttt{\_inference.\_tcp.private.ness}) create a set of SRV records from
which clients select using the RFC~2782 weighted-random algorithm, producing
load distribution and automatic failover with no scheduler required.

For Skywire-only nodes:
\begin{lstlisting}[caption={SRV record for Skywire inference endpoint.}]
KEY:   dns:_inference._skywire.batman1.private.ness
VALUE: SRV 10 100 0 02abc...def.skywire
\end{lstlisting}

Port~0 is used by convention when the transport is Skywire-native; the target
is the node's compressed public key.

\subsection{WORM Trust Chain}

Authentication does not use any DNS record type.  Instead, trust flows
through a chain of NVS WORM records that is established at node registration
time and verifiable by any client with NVS read access.

\begin{definition}[WORM trust chain]
The \emph{WORM trust chain} for an inference node is the sequence:
\[
\texttt{inf:node\_}\mathit{H} \xrightarrow{\texttt{<node>}} \texttt{node:}\mathit{name}
  \xrightarrow{\texttt{<owner>}} \texttt{id:}\mathit{username}
  \xrightarrow{\texttt{<pubkey>}} \mathit{key}
\]
Each arrow represents a typed reference field in a WORM XML document.
The chain terminates at the operator's privatenessID public key, which is
committed to the blockchain at identity registration time.
\end{definition}

The \texttt{node:} record carries both the \texttt{<owner>} reference to the
operator's \texttt{id:} record and the node's signing public key used for
TLS or Noise protocol handshake verification.  No DNS record participates
in the authentication path.

\begin{property}[Trust termination]
Trust for any inference endpoint terminates at the \texttt{id:\textit{username}}
NVS record at the tail of its WORM trust chain.  No intermediate CA, no DNS
record type, no root store, and no corporate entity participates.  Removing
or altering any record in the chain requires the controlling private key or
a Bitcoin-scale hash rate attack against the Emercoin blockchain.
\end{property}

\subsection{Five-Step Discovery Protocol}

\begin{enumerate}[label=\textbf{Step \arabic*:}, leftmargin=*, wide=0pt]
  \item \textbf{Enumerate.}  Client queries \texttt{ness:index:inf} (NVS
    prefix scan or index record) to obtain the set of known
    \texttt{inf:node\_*} keys.  Alternatively, client queries
    \texttt{dns:\_inference.\_tcp.private.ness} SRV records directly for
    transport addressing.

  \item \textbf{Select.}  Client applies the RFC~2782 weighted-random
    selection algorithm to the SRV record set, obtaining a
    \textit{(priority, weight, port, target)} tuple for transport.
    Client also identifies the corresponding \texttt{inf:node\_\textit{H}}
    key for capability lookup.

  \item \textbf{Resolve trust chain.}  Client reads
    \texttt{inf:node\_\textit{H}} WORM XML $\to$ extracts \texttt{<node>}
    reference $\to$ reads \texttt{node:\textit{name}} WORM XML $\to$
    extracts \texttt{<owner>} reference $\to$ reads
    \texttt{id:\textit{username}} WORM XML $\to$ extracts public key.
    All three reads are against Emercoin NVS directly.

  \item \textbf{Authenticate.}  Client opens a TLS or Noise protocol
    connection to the SRV-resolved endpoint.  Verifies that the node's
    presented certificate or ephemeral key matches the public key from the
    \texttt{node:} WORM record, which is itself anchored to the
    \texttt{id:} identity.  \emph{Connection is terminated immediately on
    mismatch.  No fallback is attempted.}

  \item \textbf{Negotiate.}  Client uses the \texttt{inf:node\_\textit{H}}
    WORM data (already fetched in Step~3) for model, context window, and
    pricing confirmation.  Client initiates inference session and begins
    EMC micro-payment stream.
\end{enumerate}

Figure~\ref{fig:chain} illustrates the trust chain resolved in Step~3.

\begin{figure}[ht]
\begin{center}
\begin{tabular}{l}
\texttt{inf:node\_H} $\xrightarrow{\texttt{<node>}}$ \texttt{node:batman1}\\
\texttt{node:batman1} $\xrightarrow{\texttt{<owner>}}$ \texttt{id:username}\\
\texttt{id:username} $\xrightarrow{\texttt{<pubkey>}}$ \textit{signing key}\\
\textit{signing key} $\xrightarrow{\text{verify}}$ \textit{node TLS/Noise cert}\\
\textit{cert verified} $\rightarrow$ \textit{inference traffic authorized}\\
\end{tabular}
\end{center}
\caption{WORM trust chain: \texttt{inf:} $\to$ \texttt{node:} $\to$ \texttt{id:} $\to$ public key $\to$ certificate.  No DNS record type is used for authentication.}
\label{fig:chain}
\end{figure}

\subsection{Trust Chain Enforcement}

The WORM trust chain is enforced by the PrivatenessTools client library at
Step~4.  A client that skips Step~3 connects without on-chain authentication;
this is a configuration error, not a fallback mode.  Because all three
WORM records (\texttt{inf:}, \texttt{node:}, \texttt{id:}) are read from NVS
rather than from DNS, the authentication path has no dependency on DNS record
semantics and is not affected by DNS-layer manipulation, seizure, or the
compromise of any certificate authority.

% ─── 6. ECONOMICS ────────────────────────────────────────────────────────────
\section{Proof-of-Stake Economic Model}
\label{sec:economics}

\subsection{NVS Operation Costs}

Writing a record to Emercoin NVS requires a small EMC fee proportional to
the record's key length and value size, and the desired expiry period.  A
typical \texttt{inf:} WORM XML record of approximately 800 bytes with a
one-year expiry incurs a fee on the order of 1--10~EMC at current network
parameters.  Renewal is required annually.

\subsection{Proof-of-Stake Rewards}

Emercoin operates a hybrid PoW/PoS consensus with a fixed PoS reward rate
of approximately 6\% per annum \cite{emercoin}.  Coin age maturity is 30~days:
a UTXO must remain unspent for 30 days before it becomes eligible to generate
a PoS block.  Staking rewards are proportional to the holder's balance and
accrue continuously while the wallet is online and unlocked for minting.

\begin{property}[Net-negative operational cost]
For a node operator holding an EMC staking balance $B$, the annual staking
reward is approximately $R = 0.06 \cdot B$.  For any $B$ such that
$R > C_{\mathrm{NVS}} + C_{\mathrm{renewal}}$---where $C_{\mathrm{NVS}}$
is the total cost of NVS write operations and $C_{\mathrm{renewal}}$
the total cost of annual record renewals---the net operational cost of
maintaining sovereign inference infrastructure is negative: the chain
covers costs from staking income before the first inference request is
served.
\end{property}

At current NVS fee parameters (1--10~EMC per record per year for a full
node registration comprising \texttt{id:}, \texttt{node:}, \texttt{inf:},
and \texttt{dns:} records), the staking balance required to break even
is small relative to any meaningful operational holding.  The 6\% rate
is a fixed protocol parameter, not subject to operator discretion.

\subsection{Per-Token Revenue}

Inference revenue in EMC is strictly additive.  A node generating EMC
staking rewards while simultaneously billing inference sessions compounds
both income streams.  The minimum viable economic configuration is a
small EMC stake covering operational costs with zero inference revenue;
any inference traffic is margin.

% ─── 7. HARDWARE ─────────────────────────────────────────────────────────────
\section{Commodity Hardware Deployment}
\label{sec:hardware}

A central claim of this paper is that sovereign inference is achievable on
commodity hardware.  We characterize three reproducible deployment tiers.

\subsection{Tier~0: Single-Board Computers (Pi~3 and above)}

The Raspberry Pi~3 and its successors represent the established baseline.
Mainline Linux kernel support, large community, well-understood supply chain.
Quantized 7B and 8B parameter models execute at practical token rates on
this class of hardware.  This tier establishes the \emph{conservative lower
bound}: if sovereign inference works here, it works everywhere.

\subsection{Tier~1: Directly Supported SOC TV Boxes}

Consumer Android TV hardware running commodity system-on-chip designs has
been actively brought up by the Armbian community \cite{armbian} with
mainline or near-mainline kernel support.  This hardware is designed for
mass-market consumer electronics and therefore carries aggressive price
compression not present in maker-market single-board computers.  A directly
Armbian-supported TV box running a quantized 8B model constitutes a sovereign
inference node with a bill of materials at consumer
retail pricing.

\begin{definition}[Minimum Viable Node Floor]
The \emph{minimum viable node floor} is the lowest reproducible bill-of-materials
cost for hardware capable of running a quantized 8B parameter model under
Armbian Linux with PrivatenessTools installed.  At time of writing, Tier~1
hardware in this class is available at consumer electronics retail prices,
inclusive of SOC, memory, and storage.
\end{definition}

\subsection{Tier~2: Community-Ported SOC TV Boxes}

A third class consists of TV boxes whose SOC is not directly supported by
Armbian but for which community-maintained device trees, vendor BSP fragments,
and partial mainline patches exist.  Bootstrapping sovereign inference on
Tier~2 hardware is more complex but demonstrates that the minimum viable node
floor is an \emph{active variable}, currently moving downward as the Armbian
community extends coverage.

The conservative reproducibility claim in the abstract uses Tier~0 and Tier~1
only.  Tier~2 appears as a directional indicator: the hardware floor observed
today is an upper bound on the hardware floor in 18 months.

\subsection{The 127.0.0.1 Locality Doctrine}

The architecture enforces a transport preference hierarchy:

\[\texttt{localhost} \succ \texttt{LAN node} \succ \texttt{Skywire peer} \succ \texttt{inet node}\]

A client never contacts a commercial inference API when a local or mesh-local
sovereign node is available.  This preference is not a configuration option;
it is the default client behavior, consistent with the founding philosophy
of the PrivateNess project \cite{litepaper}.

% ─── 8. SECURITY ANALYSIS ────────────────────────────────────────────────────
\section{Security Analysis}
\label{sec:security}

\subsection{Threat Model}

We consider four adversary classes:

\begin{enumerate}[label=\textbf{A\arabic*.}]
  \item \textbf{Network adversary.}  Controls BGP routing, DNS resolvers, TLS
    CAs, and internet chokepoints.  Goal: prevent access to inference services
    or redirect clients to adversarial nodes.

  \item \textbf{Platform adversary.}  Controls cloud infrastructure, payment
    processors, and app stores.  Goal: terminate inference node operation through
    service suspension or payment rail deplatforming.

  \item \textbf{Regulatory adversary.}  Issues legal demands to operators, registrars,
    or CA operators.  Goal: compel removal of service or disclosure of user identity.

  \item \textbf{Malicious node operator.}  A legitimate-appearing node operator who
    serves manipulated inference outputs.
\end{enumerate}

\subsection{Defense Against A1 (Network Adversary)}

\textbf{BGP blocking.}  Skywire transport does not use IP routing.  A Skywire-only
node has no BGP-reachable address.  BGP-level blocking is ineffective.

\textbf{DNS seizure.}  EmerDNS records are stored on the Emercoin blockchain.
There is no registrar, registry operator, or DNS hosting provider to serve with
a seizure order.  Record removal requires the operator's private key or a Bitcoin-scale
hash rate attack.

\textbf{CA compromise.}  The WORM trust chain does not involve any X.509 CA.
A CA-issued certificate for a claimed PrivateNess endpoint carries no weight
because client verification checks the NVS-anchored public key from the
\texttt{node:} WORM record directly, independent of any CA.

\textbf{Endpoint blocking.}  Inet-transport nodes are reachable only by clients
who know their SRV-resolved address.  Nodes with multiple transport declarations
fall back from inet to Skywire automatically.

\subsection{Defense Against A2 (Platform Adversary)}

Emercoin is a decentralized blockchain.  There is no hosting provider to
suspend.  EMC payment requires no payment processor.  PrivatenessTools is
open-source and distributable without an app store.  Node hardware is
consumer electronics.  There is no platform to leverage.

\subsection{Defense Against A3 (Regulatory Adversary)}

A node operator's legal identity is not required by the protocol.
The privatenessID is a cryptographic key binding with no mandatory
personal information.  A regulatory demand directed at a registrar
finds no registrar; directed at a hosting provider finds no hosting
provider; directed at a CA finds no CA.  An operator in a permissive
jurisdiction running Skywire-only transport is effectively unreachable
by foreign regulatory action.

\subsection{Defense Against A4 (Malicious Node Operator)}
\label{sec:a4}

The WORM trust chain (Steps~3--4 of the discovery
protocol) ensures that a client connects only to a node whose signing key is
committed on-chain under the declared operator identity.  A malicious
node \emph{that has registered a legitimate \texttt{id:} record} can
serve manipulated outputs; this is addressed by the reputation layer
(\texttt{rep:} prefix, Section~\ref{sec:rep}) which is an open problem.

A malicious node that has \emph{not} registered a legitimate identity
cannot pass Step~5 authentication and is unreachable by compliant clients.

\subsection{Entropy as a Security Primitive}
\label{sec:entropy-threat}

\begin{property}[Operator entropy control]
A centralized inference service whose operator controls the sampling RNG
can bias, watermark, or suppress token selections in ways that are
statistically undetectable to individual users and require population-scale
analysis to identify.
\end{property}

This attack surface is implicit in every commercial inference deployment
but has not been previously stated as a formal threat.  Temperature-based
sampling selects tokens proportional to $\exp(l_i / T)$ where $l_i$ are
logits and $T$ is temperature.  An operator who controls the RNG can
consistently bias the sampling distribution without altering any visible
parameter.

The PrivateNess stack uses UHEPRNG \cite{pyuheprng}, a user-controlled
entropy source, as the sampling RNG.  Because the entropy is generated
by the client or by the sovereign node's own PRNG service---not by an
operator-controlled backend---the attack surface is closed structurally.

\subsection{Fiber Hostile Node Resistance}

The Skywire Fiber consensus model \cite{skywire} guarantees safety with
a single honest node.  The failure threshold for the PrivateNess inference
network is therefore $n - 1$ of $n$ nodes, not a majority fraction.  The
last honest node in the network constitutes a valid, operational inference
service.  This is the lower bound on network availability; the expected
operating point is far above it.

% ─── 9. LIMITATIONS ──────────────────────────────────────────────────────────
\section{Limitations and Open Problems}
\label{sec:limitations}

\subsection{Provider Reputation Protocol}
\label{sec:rep}

The \texttt{rep:} prefix is specified (Table~\ref{tab:prefixes}) but not yet
implemented.  Attestation records require a trust weighting scheme that
does not recurse to a central authority.  We propose weighting attestations
by the anchor age of the attesting \texttt{id:} record: older anchors cost
more to establish and are therefore harder to Sybil-attack.  The AuxPoW
cost of registering a new \texttt{id:} record provides a baseline
Sybil-resistance floor.  A complete reputation protocol is left to future
work.

\subsection{Inference Latency}

Skywire mesh routing adds per-hop latency.  Quantized 8B models on Tier~1
and Tier~2 hardware generate tokens at rates of 5--20 tokens per second
depending on context length.  The combined latency profile makes real-time
interactive use marginal on current hardware; batch generation, document
analysis, and asynchronous workflows are the practical deployment targets
in the near term.  This is a positioning constraint, not a protocol deficiency.

\subsection{Liveness Signalling Without Record Mutation}

A node that goes offline does not automatically expire its SRV or
\texttt{inf:} records.  Clients discover stale records and incur failed
connection attempts before falling back to an alternative SRV entry.
A liveness layer is required; candidates include:
(a)~proof-of-liveness WORM records with recent timestamp updates,
(b)~Skywire transport failure as an implicit liveness signal, or
(c)~a gossip protocol over Skywire carrying liveness beacons.
None of these is implemented.

\subsection{NVS Block Time and Real-Time Load Balancing}

Emercoin targets a block interval of approximately 10~minutes (the observed
network average is around 8--9~minutes), identical to Bitcoin's cadence
\cite{emercoin}.  NVS-based service discovery records are therefore updated
on a 10-minute timescale at best.  This makes NVS unsuitable for sub-minute
dynamic load balancing.  The SRV priority/weight mechanism handles
coarse-grained balancing at registration time.  Fine-grained real-time
routing requires a second-layer gossip protocol over Skywire.
This is an open engineering problem.

\subsection{Trust Chain Enforcement Dependency on Client Library}

WORM trust chain verification is not automatic; it requires clients to use
the PrivatenessTools library or a compatible implementation that performs
the full \texttt{inf:} $\to$ \texttt{node:} $\to$ \texttt{id:} resolution
and key verification at Step~4.  A client that skips this sequence connects
without on-chain authentication.  This is an implementation dependency, not
a protocol flaw, but it creates a correctness requirement on client software
that is not self-enforcing.

% ─── 10. RELATED WORK ────────────────────────────────────────────────────────
\section{Related Work}
\label{sec:related}

\textbf{Decentralized AI platforms.}  Bittensor \cite{bittensor} creates a
token-incentivized market for AI inference and training, but discovery and
authentication remain mediated by a centralized validator set.  Akash Network
\cite{akash} provides decentralized compute but does not address AI-specific
concerns such as entropy control or prompt privacy.  Neither system terminates
trust on a blockchain-anchored identity without intermediate infrastructure.

\textbf{Federated inference.}  Federated learning literature (e.g.,
\cite{federated}) addresses training data privacy but does not address inference
endpoint trust, service discovery, or payment.

\textbf{DNS-based service discovery.}  RFC~2782 SRV records are widely used
for VOIP \cite{sip}, XMPP \cite{xmpp}, and other distributed services.
Our contribution is the application of SRV to AI inference service discovery
with blockchain-backed records, where SRV handles transport addressing only
and authentication is handled entirely within the NVS WORM record chain.

\textbf{Self-sovereign identity.}  The W3C DID specification \cite{w3cdid}
addresses decentralized identity but does not specify a complete service
discovery, authentication, and payment stack.  The privatenessID system
predates the DID specification's finalization and is optimized for the
Emercoin NVS substrate rather than a generic resolver architecture.

% ─── 11. CONCLUSION ──────────────────────────────────────────────────────────
\section{Conclusion}
\label{sec:conclusion}

We have presented a complete architecture for sovereign inference-as-a-service:
AI inference that is discovered via blockchain-hosted SRV records, authenticated
via a WORM trust chain rooted entirely in on-chain NVS identity records,
transported over a public-key-routed mesh network, paid in a decentralized
currency, and operated on commodity hardware available at consumer electronics retail prices.

The five contributions are independent but mutually reinforcing.  The
economic model makes node operation self-sustaining.  The hardware floor
makes participation accessible.  The SRV + WORM trust chain makes the
network discoverable with standard DNS tooling while keeping authentication
entirely on-chain.  The privatenessID record chain makes authentication
independent of any certificate authority, DNS record type repurposing,
or external trust infrastructure.
The sovereign entropy source closes an attack surface that is invisible
in the centralized inference model.

Three open problems remain: provider reputation, liveness signalling, and
real-time load balancing.  Each is tractable given the existing infrastructure;
none requires changes to the core protocol.

The minimum viable sovereign inference node is not a research artifact.
It runs today, on hardware you can buy at a consumer electronics retailer,
on a network whose trust roots were committed to the blockchain in 2013 and
2020.  Decentralized AI infrastructure, operated at the edge, at commodity
cost, by the people who use it, is the structural answer to centralized
inference dependency.

% ─── ACKNOWLEDGEMENTS ────────────────────────────────────────────────────────
\section*{Acknowledgements}

The author thanks Aleksej Sokolov (Alex014) for the PrivatenessTools
implementation and WORM XML format; Oleg Khovayko for the Emercoin NVS
infrastructure and hack51 AuxPoW security framework; and Brandon Smietana
and Moses Narrow for the Skywire mesh network and Fiber consensus model.
Prior art is anchored on-chain at Emercoin blocks \#445669, \#780915,
and \#780931.

% ─── REFERENCES ──────────────────────────────────────────────────────────────
\begin{thebibliography}{99}

\bibitem{emercoin}
Emercoin Development Team.
\textit{Emercoin: Blockchain with Name-Value Storage}.
\url{https://emercoin.com}, 2013.

\bibitem{hack51}
O.~Khovayko.
\textit{hack51: AuxPoW Security Coefficient Framework for Emercoin NVS}.
Internal technical specification, Emercoin project.

\bibitem{privatenesstools}
A.~Sokolov (Alex014).
\textit{PrivatenessTools: Console Tools for the PrivateNess Network}.
\url{https://github.com/NESS-Network/PrivatenessTools}, 2020--2025.

\bibitem{pyuheprng}
A.~Sokolov (Alex014).
\textit{pyuheprng: Uniform Hash-based Entropy PRNG}.
\url{https://github.com/Alex014/pyuheprng}.

\bibitem{identitybedrock}
PrivateNess Network.
\textit{Identity Bedrock: Rompel-Grounded Sovereign Identity with
Ed25519/ML-DSA Blockchain Anchoring and Physically-Bound Authentication}.
PrivateNess Network technical paper, 2024--2025.

\bibitem{skywire}
B.~Smietana and M.~Narrow.
\textit{Skywire: A Peer-to-Peer Networking Platform}.
Skycoin project, \url{https://skywire.skycoin.com}.

\bibitem{litepaper}
PrivateNess Network.
\textit{PrivateNess Network Litepaper}.
\url{https://privateness.network}, October~9, 2020.

\bibitem{armbian}
Armbian Community.
\textit{Armbian: Linux for ARM Development Boards}.
\url{https://armbian.com}.

\bibitem{ipfs}
J.~Benet.
\textit{IPFS --- Content Addressed, Versioned, P2P File System}.
arXiv:1407.3561, 2014.

\bibitem{ethereum}
G.~Wood.
\textit{Ethereum: A Secure Decentralised Generalised Transaction Ledger}.
Ethereum Foundation, 2014.

\bibitem{rfc2782}
A.~Gulbrandsen, P.~Vixie, and L.~Esibov.
\textit{A DNS RR for Specifying the Location of Services (DNS SRV)}.
RFC~2782, IETF, February~2000.

\bibitem{bittensor}
J.~Fishman and A.~Psoras.
\textit{Bittensor: A Peer-to-Peer Intelligence Market}.
\url{https://bittensor.com}, 2021.

\bibitem{akash}
G.~Adams et~al.
\textit{Akash Network: Decentralized Cloud Computing}.
\url{https://akash.network}, 2020.

\bibitem{federated}
B.~McMahan et~al.
\textit{Communication-Efficient Learning of Deep Networks from Decentralized Data}.
Proceedings of AISTATS, 2017.

\bibitem{sip}
J.~Rosenberg et~al.
\textit{SIP: Session Initiation Protocol}.
RFC~3261, IETF, June~2002.

\bibitem{xmpp}
P.~Saint-Andre.
\textit{Extensible Messaging and Presence Protocol (XMPP): Core}.
RFC~6120, IETF, March~2011.

\bibitem{w3cdid}
M.~Sporny et~al.
\textit{Decentralized Identifiers (DIDs) v1.0}.
W3C Recommendation, July~2022.

\end{thebibliography}

\end{document}
